\SourceLeaf{291}%
\[
\begin{aligned}
\sin(Q+q)&=\sin Q\cos\tfrac{1}{2}(h-\delta)+\cos Q\sin\tfrac{1}{2}(h-\delta)\cos\chi\\
\sin(Q-q)&=\sin Q\cos\tfrac{1}{2}(h+\delta)-\cos Q\sin\tfrac{1}{2}(h+\delta)\cos\chi
\end{aligned}
\]
und ihre Summe, geteilt durch \(2\cos Q\), ergibt
\[
\operatorname{tang}Q\mathord{.}(\cos q-\cos\tfrac{1}{2}h\mathord{.}\cos\tfrac{1}{2}\delta)=-\cos\tfrac{1}{2}h\sin\tfrac{1}{2}\delta\cos\chi
\]
\indent Da nun \(\cos q-\cos\tfrac{1}{2}h\mathord{.}\cos\tfrac{1}{2}\delta\) offenbar eine Größe zweiter Ordnung ist, sind auch \(\sin\tfrac{1}{2}\delta\cos\chi\) und \(\delta\cos\chi\) von dieser Ordnung. Da \(\lambda\) den Faktor \(\cos\chi^2\) enthält, ist folglich \(\lambda hh\delta\) von vierter und \(\lambda h\delta\delta\) von fünfter Ordnung. Damit ist es gerechtfertigt, diese Glieder wegzulassen.

Unter der angegebenen Voraussetzung ist das Endergebnis dieser Entwicklung somit in den folgenden Formeln enthalten. Darin stehen anstelle der Tangenten von \(\psi^0,\psi'\) die Bögen selbst
\[
\begin{aligned}
\psi^0&=-\frac{ee\cos P\sin P\sin\chi\cos\chi^2h^3}{12\cos\varphi\cos\theta}\\
\psi'&=+\frac{ee\cos P\sin P\sin\chi\cos\chi^2h^3}{12\cos\varphi\cos\theta}
\end{aligned}
\]
\Article{16}
Die Berechnung des Dreiecksnetzes auf der Kugel gliedert sich in drei Hauptteile:
\begin{list}{}{\setlength{\leftmargin}{1.8em}\setlength{\labelwidth}{1.4em}\setlength{\labelsep}{.4em}\setlength{\itemsep}{0pt}}
\item[1)] die Ausgleichung der Winkel nach allen Bedingungsgleichungen, die sich aus der Beschaffenheit des Netzes ergeben.
\item[2)] die Berechnung sämtlicher Dreiecksseiten.
\item[3)] \begingroup\rightskip=0pt plus 1em
die Bestimmung der Längen und Breiten der Dreieckspunkte zusammen mit der Orientierung der Dreiecksseiten, die von jedem dieser Punkte ausgehen.\par\endgroup
\end{list}
\hspace*{\parindent}Die Längen und Breiten auf der Kugel werden dann in die wahren Längen und Breiten auf dem Sphäroid umgerechnet. Die Längen werden durch den konstanten Divisor \(\alpha\) geteilt. Für die Breiten verwendet man die hier beigefügte Hilfstafel oder eine andere, eigens auf ähnliche Weise erstellte Tafel, wenn es einen Grund gibt, einen anderen Normalparallelkreis zu wählen.

Die beiden ersten, auf bekannten Grundlagen beruhenden Arbeiten übergehe ich und füge hier noch einiges zur dritten hinzu. Diese läuft auf die Lösung der folgenden Aufgabe hinaus\footnote{Da diese Aufgabe hier für sich betrachtet wird, können einige Buchstaben hier ohne Nachteil in einer anderen Bedeutung als oben verwendet werden.}: Aus der in Bogenteilen ausgedrückten
